\documentclass[lineno]{JFM-FLM_Au}

\usepackage{xcolor}
\usepackage{tikz}
\usetikzlibrary{calc,arrows.meta}
\usepackage{xcolor}  % if not already loaded
\usepackage{amsthm}
\newcommand{\tbd}[1]{\textcolor{red}{\ensuremath{\langle\mbox{\small #1}\rangle}}}
\definecolor{jfmblue}{RGB}{0,56,168}
\def\xshift{1.5}
\hypersetup{
  colorlinks=true,
  citecolor=blue,
  linkcolor=jfmblue,
  urlcolor=jfmblue
}

\newtheorem{lemma}{Lemma}
\newtheorem{corollary}{Corollary}

\lefttitle{Sharma et al.}
\righttitle{Journal of Fluid Mechanics}

\title{}

% % \author{Sparsh Sharma\aff{1}, Author 2\aff{1}, Author 3\aff{1},  \textcolor{red}{(Stéphane Moreau?)}\aff{2}
%  \and Michaela Herr\aff{1}}

\affiliation{\aff{1}German Aerospace Center (DLR), 38108 Braunschweig, Germany
\aff{2}Université de Sherbrooke, Sherbrooke, J1K 2R1, Canada}

\corresau{Sparsh Sharma, sparsh.sharma@dlr.de}

\begin{document}
\maketitle

\begin{abstract}
This file contains information for authors planning to submit a paper to the {\it Journal of Fluid Mechanics}. The document was generated in {\LaTeX} using the JFM class file and supporting files provided on the JFM website \href{https://www.cambridge.org/core/journals/journal-of-fluid-mechanics/information/author-instructions/preparing-your-materials}{here}, and the source files can be used as a template for submissions (please note that this is mandatory for {\it JFM Rapids}). Full author instructions can be found on the \href{https://www.cambridge.org/core/journals/journal-of-fluid-mechanics/information/author-instructions}{JFM website}. The present paragraph appears in the \verb}abstract} environment. All papers should feature a single-paragraph abstract of no more than 250 words which must not spill onto the second page of the manuscript.
\end{abstract}

\begin{keywords}
Authors should not enter keywords on the manuscript, as these must be chosen by the author during the online submission process and will then be added during the typesetting process (see \href{https://www.cambridge.org/core/services/aop-file-manager/file/61436b61ff7f3cfab749ce3a/JFM-Keywords-Sept-2021.pdf.}{Keyword PDF} for the full list).  Other classifications will be added at the same time.
\end{keywords}

%{\bf MSC Codes }  {\it(Optional)} Please enter your MSC Codes here

% =========================================================================
% Introduction — Turbulence distortion and the pre-impact spectral tensor
% in airfoil leading-edge noise
% Target journal: Journal of Fluid Mechanics (natbib, \citep / \citet)
% =========================================================================

% \section{Introduction}
% \label{sec:introduction}

% The interaction of free-stream turbulence with the leading edge of a lifting
% surface is among the dominant broadband noise mechanisms in low-Mach-number
% aeroacoustics, governing the acoustic signature of fans, propellers,
% turbomachinery stages operating in wakes or distorted inflow, and wind
% turbines in the atmospheric boundary layer
% \citep{Amiet1975,PatersonAmiet1976,Moriarty2005,DevenportStaubsGlegg2010}.
% As trailing-edge and tonal sources have been progressively attenuated through
% geometric optimisation and flow control, turbulence-interaction noise has
% become an increasingly binding constraint on low-noise design, and the demand
% for predictive — rather than calibrated — models of this source has grown
% correspondingly \citep{KimHaeriJoseph2016,SharmaSarradjSchmidt2020,SharmaHerr2024}.

% The classical description of the generation mechanism is well established.
% An incoming vortical disturbance, decomposed into sinusoidal gusts, induces
% an unsteady loading on the airfoil whose linear transfer function was first
% derived for incompressible flow by \citet{Sears1941} and generalised to
% oblique gusts and subsonic compressible flow by \citet{Graham1970}. The
% unsteady surface pressure radiates to the far field by a scattering process
% whose theoretical foundations lie in the acoustic analogy of
% \citet{Lighthill1952}, its extension to flows containing stationary and
% moving surfaces \citep{Curle1955,FfowcsWilliamsHawkings1969}, and the
% analysis of scattering by sharp edges \citep{FfowcsWilliamsHall1970}; unified
% treatments are given by \citet{Goldstein1976} and, in the vortex-sound
% framework, by \citet{Howe2003}. Building on these elements,
% \citet{Amiet1975} formulated the canonical leading-edge noise model: the
% far-field pressure spectral density is expressed as the product of an
% aerodynamic--acoustic transfer function, derived for a semi-infinite flat
% plate, and the wavenumber spectrum of the incident velocity fluctuations.
% The accompanying experiments of \citet{PatersonAmiet1976} established the
% quantitative credibility of this formulation, and subsequent analytical
% extensions have refined the finite-chord and back-scattering corrections
% while retaining its overall architecture
% \citep{RogerMoreau2005,MoreauRoger2009}.

% The predictive power of this architecture rests on a specific set of
% assumptions about the incident turbulence. The incoming field is taken to be
% frozen in the sense of \citet{Taylor1938}, convecting toward the leading edge
% at the free-stream velocity without modification by the mean flow; it is
% usually modelled as homogeneous and isotropic, so that the full spectral
% tensor is parametrised by an intensity and an integral length scale; and, of
% the nine tensor components, only the wall-normal (upwash) velocity spectrum
% enters the prediction, since for a flat plate at zero incidence only the
% upwash produces unsteady loading \citep{Amiet1975,Goldstein1976}. Under
% these assumptions the turbulence appears in the model purely as a prescribed
% boundary datum: the entire physics of the interaction is delegated to the
% acoustic transfer function, and the spectrum evaluated far upstream is
% implicitly identified with the spectrum that forces the leading edge.

% A substantial body of experimental and computational work has tested this
% framework and delineated where it fails. \citet{PatersonAmiet1976} already
% noted an overprediction of high-frequency noise for airfoils of finite
% thickness, a trend confirmed systematically by \citet{OlsenWagner1982} and
% quantified in later studies as a thickness- and frequency-dependent
% attenuation relative to the flat-plate result
% \citep{Gershfeld2004,Lysak2013,Gill2013}. Measurements of unsteady surface
% pressure beneath grid turbulence showed that mean loading and geometry
% modify the pressure response in ways not captured by thin-airfoil theory
% \citep{MishDevenport2006}, and acoustic measurements on families of real
% airfoils demonstrated that the high-frequency noise reduction is not
% controlled by a single geometric parameter such as the leading-edge radius or
% maximum thickness \citep{DevenportStaubsGlegg2010,Hall2011}. On the
% computational side, gust--airfoil interaction in realistic mean flows
% \citep{LockardMorris1998}, high-fidelity simulations of
% aerofoil--turbulence interaction \citep{KimHaeriJoseph2016}, and, more
% recently, lattice Boltzmann simulations of grid-generated turbulence
% impinging on a NACA\,0012 airfoil
% \citep{Trascinelli2023,Piccolo2024} have resolved both the hydrodynamic
% near field and the acoustic radiation, while dedicated aeroacoustic
% facilities have characterised the inflow turbulence with increasing care
% \citep{Bowen2022}. These studies agree that finite-thickness effects are
% real, scale-dependent and acoustically significant.

% What they leave less settled is where in the causal chain those effects
% should be attributed. Most corrections developed to reconcile flat-plate
% theory with thick-airfoil data act at the level of the acoustic response:
% streamline-based modifications of the incident gust along the mean-flow
% trajectory \citep{Guidati1997,Moriarty2005}, exponential high-frequency
% attenuation factors \citep{Gershfeld2004,Lysak2013}, or panel-method
% generalisations of the transfer function to realistic geometries
% \citep{GleggDevenport2010Panel}. Each of these improves agreement within its
% calibration range, but each also absorbs the underlying physics — the
% modification of the turbulence itself before it reaches the surface — into
% an effective acoustic operator. The causal chain implied by the flow is
% longer than the one represented in the models: the upstream turbulence
% tensor is first transformed by the decelerating, straining flow ahead of the
% stagnation region, and only the resulting pre-impact state is available to
% force the unsteady loading. Classical leading-edge noise theories describe
% the scattering and acoustic response once an incoming gust or spectrum has
% been specified; comparatively little attention has been given to the
% transformation that converts the measured upstream turbulence into the
% effective state experienced by the leading edge.

% The physical and mathematical language appropriate to this earlier stage
% exists, and is that of rapid distortion theory (RDT). Following the analysis
% of a single Fourier mode under uniform straining by \citet{Taylor1935},
% \citet{BatchelorProudman1954} derived the linear response of homogeneous
% turbulence to rapid mean deformation, and \citet{Hunt1973} extended the
% theory to flow around two-dimensional bluff bodies, identifying the distinct
% roles of mean-strain distortion, inviscid blocking by the surface, and the
% concentration of vorticity near the stagnation streamline; the blocking
% mechanism near a plane boundary was analysed separately by
% \citet{HuntGraham1978}. Hot-wire measurements upstream of bluff bodies
% confirmed the predicted behaviour, with large-scale wall-normal fluctuations
% suppressed by blocking and smaller scales amplified by straining, the
% crossover being set by the ratio of eddy size to body dimension
% \citep{Bearman1972,BritterHuntMumford1979}. In the aeroacoustic context,
% \citet{GoldsteinAtassi1976} showed that the unsteady airfoil response is
% fundamentally altered when the distortion of the incident gust by the steady
% potential flow is retained, and \citet{Goldstein1978} generalised the
% splitting of the disturbance field to arbitrary obstacles;
% \citet{HuntCarruthers1990} review the scope and limitations of the linear
% theory. These results establish that the quantity transformed by the
% leading-edge flow is not a scalar intensity but the full spectral tensor:
% the deformation acts differently on different components and different
% wavenumbers, generating anisotropy that is naturally described by the
% Reynolds-stress anisotropy invariants of \citet{LumleyNewman1977}, and
% modified further by pressure--strain redistribution beyond the reach of
% linear theory \citep{SpezialeSarkarGatski1991,Pope2000}. Recent work has
% begun to import these ideas into noise prediction: RDT-corrected upwash
% spectra improve Amiet-type predictions for thick airfoils
% \citep{deSantana2016,dosSantos2023}, distortion has been invoked to explain
% the noise reduction of porous leading edges \citep{Zamponi2020}, and
% combined experimental--numerical studies have measured the one-dimensional
% spectral modification directly in the stagnation region \citep{Piccolo2024}.

% Taken together, however, the literature characterises this pre-impact
% transformation only in reduced form — as a correction to a single velocity
% spectrum, a streamline-integrated velocity ratio, or a fitted attenuation
% function. Several questions of principle remain open. It has not been
% established empirically whether the mapping from the upstream spectral
% tensor $\Phi^{\infty}_{ij}$ to the pre-impact tensor $\Phi^{\mathrm{LE}}_{ij}$
% can be represented by a scalar, frequency-dependent gain, or whether it is
% irreducibly tensorial, coupling components and redistributing energy
% anisotropically. The scale dependence of the mapping — in particular,
% whether the leading-edge radius $r_{\mathrm{LE}}$ acts as the geometric
% filter separating strongly and weakly distorted scales through the
% dimensionless wavenumber $\kappa = k\,r_{\mathrm{LE}}$ — has been inferred
% from acoustic far-field trends \citep{Gill2013,DevenportStaubsGlegg2010} but
% not measured directly on the turbulence field surrounding the leading edge.
% Nor is it known to what extent the transformation departs from linear RDT,
% through pressure--strain redistribution and finite-time nonlinear effects,
% in the regime of practical interest. Until these questions are answered,
% the common practice of inserting an upstream-measured spectrum directly into
% a leading-edge noise model rests on an assumption whose error is neither
% quantified nor even of known sign across frequency.

% The present paper addresses these questions by isolating the pre-impact
% transformation itself as the object of study. We perform wall-modelled
% large-eddy simulations, using a lattice Boltzmann method, of grid-generated
% turbulence convecting onto a NACA\,0012 airfoil at chord-based Reynolds
% number $Re_c \approx 5.1\times10^{5}$ and free-stream Mach number
% $M = 0.059$, and we compare the turbulence state on planar reference patches
% upstream of the airfoil with that on a toroidal sampling surface enclosing
% the leading edge. The comparison is organised around the hypothesis of a
% distortion mapping,
% $\Phi^{\mathrm{LE}}_{ij}(\boldsymbol{k},\omega)
%  = \mathcal{D}_{ijmn}(\boldsymbol{k},\theta)\,
%   \Phi^{\infty}_{mn}(\boldsymbol{k},\omega)$,
% in which a fourth-order, scale- and orientation-dependent operator
% $\mathcal{D}$, structurally motivated by rapid distortion theory, relates
% the two spectral tensors. We emphasise that $\mathcal{D}$ is introduced
% here as a diagnostic and organising framework rather than as an established
% closed-form model: the simulations provide its diagonal and azimuthally
% projected components through component-wise spectral gains $G_q(f)$ and
% local-frame distortion maps, together with Reynolds-stress budgets and
% Lumley-invariant trajectories that constrain its tensorial structure. A
% preliminary account of this approach was given in
% \citet{Sharma2026}. The results demonstrate that the transformation is
% scale-dependent, with the transition between strong and weak distortion
% organised by $\kappa = k\,r_{\mathrm{LE}}$; that it is component-dependent
% and therefore not reducible to a scalar correction; and that the invariant
% trajectories exhibit a transient eigenvalue reordering inconsistent with
% purely linear distortion, implicating pressure--strain redistribution in the
% pre-impact dynamics.

% The remainder of the paper is organised as follows.
% Section~\ref{sec:framework} develops the mathematical framework that casts
% turbulence distortion as an operator mapping between spectral tensors and
% identifies $\kappa$ and the integrated-strain parameter $\theta$ as the
% governing dimensionless groups. Sections~\ref{sec:method}
% and~\ref{sec:setup} describe the lattice Boltzmann solver, the subgrid-scale
% and wall modelling, and the computational configuration, including the
% turbulence-generating grid and the two-surface sampling strategy.
% Sections~\ref{sec:flowfield}--\ref{sec:invariants} characterise the
% evolution of the mean flow, Reynolds stresses and anisotropy invariants
% between the upstream reference and the leading edge.
% Section~\ref{sec:spectra} presents the patch-based spectral comparison, the
% spectral gains $G_q(f)$ and the local-frame distortion maps that constitute
% the empirical characterisation of $\mathcal{D}$. Section~\ref{sec:conclusions}
% summarises the findings and discusses their implications for the structure
% of predictive leading-edge noise models.

% % Standalone-insertable figure for the JFM manuscript.
% % Requires in preamble: \usepackage{tikz,graphicx,amsmath}
% %                       \usetikzlibrary{calc,arrows.meta}
% % (all already present in your new-aiaa preamble)

% % \begin{figure}
% % \centering
% % \begin{tikzpicture}[>=Latex, font=\small]
% %   % ---- mean flow ----
% %   \foreach \y in {-1.3,-0.65,0.65,1.3}
% %     \draw[->,gray] (-5.0,\y) -- (-4.3,\y);
% %   \node[gray] at (-4.70,1.98) {$U_\infty\boldsymbol{e}_x$};
% %   % ---- upstream reference plane ----
% %   \draw[dashed,thick] (-3.6,-1.9) -- (-3.6,2.1);
% %   \node at (-3.6,2.40) {$\mathcal{S}_\infty:\;\Phi^{\infty}_{ij}$};
% %   % ---- airfoil: nose circle r_LE = 0.85 centred at origin, tail (6.5,0) --
% %   \draw[thick,fill=gray!12]
% %     (0,0.85) .. controls (2.3,0.60) and (5.0,0.20) .. (6.5,0)
% %              .. controls (5.0,-0.20) and (2.3,-0.60) .. (0,-0.85)
% %     arc (-90:-270:0.85);
% %   \draw[densely dotted] (0,0) circle (0.85);    % osculating nose circle
% %   \node[circle,fill,inner sep=0.7pt] at (0,0) {};
% %   \draw[->] (0,0) -- (35:0.85);
% %   \node at (1.12,0.55) {$r_{\mathrm{LE}}$};
% %   % ---- pre-impact surface (conformal, offset delta_n) ----
% %   \draw[dashed,thick] (95:1.30) arc (95:265:1.30);
% %   \draw[gray!70] (243:1.30) -- (-1.45,-1.55);
% %   \node[anchor=east] at (-1.40,-1.62) {$\mathcal{S}_{\mathrm{LE}}:\;\Phi^{\mathrm{LE}}_{ij}$};
% %   \draw[<->] (-0.88,-0.16) -- (-1.27,-0.16);
% %   \node[below] at (-1.07,-0.20) {$\delta_n$};
% %   % ---- stagnation line, azimuth ----
% %   \draw[gray] (-3.6,0) -- (-1.35,0);
% %   \node[gray,below] at (-2.55,0) {$\theta=0$};
% %   \draw[->] (-1.58,0) arc (180:123:1.58);
% %   \node at (-1.78,0.72) {$\theta$};
% %   % ---- local frame at drawing angle 120 deg (theta = 60 deg) ----
% %   \coordinate (P) at (120:1.30);
% %   \node[circle,fill,inner sep=0.7pt] at (P) {};
% %   \draw[->,thick] (P) -- ++(120:0.60);
% %   \node[anchor=south,inner sep=2pt] at ($(P)+(120:0.60)$)
% %         {$\boldsymbol{e}_r=\boldsymbol{e}_n$};
% %   \draw[->,thick] (P) -- ++(30:0.60);
% %   \node[anchor=west,inner sep=2pt] at ($(P)+(30:0.60)$)
% %         {$\boldsymbol{e}_\theta=\boldsymbol{e}_t$};
% %   % ---- eddy glyphs: sphere -> tangentially elongated ellipse ----
% %   \draw[gray!70] (-3.0,0.85) .. controls (-2.3,0.90) and (-1.9,1.05) .. (-1.55,1.35);
% %   \draw (-3.0,0.85) circle (0.20);
% %   \draw[rotate around={38:(-1.55,1.35)}] (-1.55,1.35) ellipse (0.33 and 0.11);
% %   \node[align=center] at (-2.98,1.60)
% %         {\footnotesize tangential extension,\\[-3pt]\footnotesize normal compression};
% %   % ---- dimensions / plane note ----
% %   \draw[<->] (-3.6,-1.95) -- (-0.85,-1.95);
% %   \node[below] at (-2.22,-1.95) {$\ell_\infty$};
% %   \node[align=left] at (4.7,1.55)
% %         {\footnotesize midspan $(x,y)$ plane;\\[-2pt]
% %          \footnotesize $\boldsymbol{e}_z$ out of plane};
% % \end{tikzpicture}
% % \caption{Configuration and definitions of \S\,\ref{sec:framework_states}.
% % Turbulence carried by the uniform stream is characterised on the upstream
% % reference plane $\mathcal{S}_\infty$ (a distance $\ell_\infty$ ahead of the
% % leading edge) and again on the conformal pre-impact surface
% % $\mathcal{S}_{\mathrm{LE}}$, offset $\delta_n \ll r_{\mathrm{LE}}$ from the
% % nose. The azimuth $\theta$ is measured from the upstream stagnation line
% % about the centre of the osculating nose circle (dotted); on that circle
% % the cylindrical and surface-aligned frames coincide,
% % $\boldsymbol{e}_r = \boldsymbol{e}_n$,
% % $\boldsymbol{e}_\theta = \boldsymbol{e}_t$. The eddy glyphs indicate the
% % sense of the strain (\ref{eq:strain_local}): tangential extension, normal
% % compression.}
% % \label{fig:framework_config}
% % \end{figure}

% % Standalone-insertable figure for the JFM manuscript.
% % Requires in preamble: \usepackage{tikz,graphicx,amsmath}
% %                       \usetikzlibrary{calc,arrows.meta}
% % (all already present in your new-aiaa preamble)



% % Standalone-insertable figure for the JFM manuscript.
% % Requires in preamble: \usepackage{tikz,graphicx,amsmath}
% %                       \usetikzlibrary{calc,arrows.meta}
% % (all already present in your new-aiaa preamble)

% % =========================================================================
% % Section 2 (EXPANDED VERSION) — Theoretical framework: pre-impact
% % distortion as a tensorial mapping.
% % Target: Journal of Fluid Mechanics (natbib, \citep / \citet).
% %
% % NOTATION (must propagate through the manuscript):
% %   \vartheta : integrated-strain (distortion) parameter  [was θ in conf. §II]
% %   \theta    : azimuthal coordinate around the leading edge, in the
% %               MIDSPAN (x,y) plane, θ = 0 on the upstream stagnation line.
% %   Local cylindrical projection acts on (u,v) in the midspan plane, w
% %   unchanged (published-PDF convention). >>> DECISION POINT: if the
% %   post-processing projected (v,w) in the (y,z) plane instead, §2.6 and
% %   eq. (trace) must be changed to match. <<<
% %
% % All citations resolve to entries in references.bib +
% % references_framework_additions.bib (all verified). No new references.
% % =========================================================================
% =====================================================================
%  Introduction section  --  strain-coupled / RDT-consistent ODT
%  Target journal: Journal of Fluid Mechanics (jfm.cls)
%
%  Preamble requirements (in your main .tex):
%     \documentclass{jfm}
%     \usepackage{amsmath,amssymb}
%  Bibliography: \bibliographystyle{jfm}  +  \bibliography{references}
%  Citation commands used: \citet (textual) and \citep (parenthetical),
%  which are provided by jfm.cls / natbib.
%
%  NOTE ON CITATIONS: every reference used below is defined in
%  references.bib and has been bibliographically verified, except where
%  a TODO comment in references.bib flags a field to confirm.
% =====================================================================

\section{Introduction}\label{sec:intro}

Turbulence--airfoil interaction noise, radiated when a lifting surface
encounters oncoming turbulence, is a dominant broadband source in axial
fans, contra-rotating rotors, wind-turbine and propeller blades, and
airframe components operating in disturbed inflow. When the integral
length scale of the incident turbulence is large compared with the
leading-edge radius, the radiated field is produced predominantly at the
leading edge, where the unsteady upwash induces a fluctuating surface
loading that scatters into sound as a non-compact dipole. Predicting this
leading-edge noise both reliably and at a cost compatible with design
iteration remains an open problem, and is the practical motivation for the
present work.

The prevailing predictive framework is the analytical theory of
\citet{Amiet1975}, which combines the linearised gust-response solution of
\citet{Sears1941} with the surface-source formulation of \citet{Curle1955}
and can be embedded within the general moving-surface analogy of
\citet{FfowcsWilliams1969}. In Amiet's theory the far-field acoustic
power spectral density is obtained as a linear functional of the
wavenumber--frequency spectrum of the upwash (airfoil-normal) velocity
component of the incoming turbulence, convolved with an airfoil response
function and an edge-scattering kernel and weighted by the spanwise
correlation of the loading. The framework has been extended to finite
chord, finite span and realistic geometry and validated against
grid-turbulence experiments \citep{PatersonAmiet1977,RogerMoreau2010},
but the structure of the prediction is unchanged: \emph{the upwash
spectrum that enters the theory is prescribed, not computed}. In practice
it is taken from an idealised homogeneous, isotropic model spectrum---most
commonly the von~K\'arm\'an form \citep{vonKarman1948}---fixed by the
upstream turbulence intensity and integral length scale and frozen in
space through Taylor's hypothesis. Because the radiated spectrum is linear
in this input, the fidelity of any leading-edge-noise prediction is
limited, first and foremost, by the fidelity of the prescribed upwash
spectrum.

The assumption that the turbulence reaching the leading edge is identical
to the undisturbed upstream turbulence is, however, physically
questionable. As the flow approaches the stagnation region the mean
velocity gradient distorts the incident turbulence over a time that, for
the energy-containing scales, can be comparable to or shorter than the
eddy-turnover time. This is the rapid-distortion regime first analysed by
\citet{BatchelorProudman1954} and developed for flow around obstacles by
\citet{Hunt1973} and \citet{Goldstein1978}: the turbulence becomes
anisotropic and its spectral content is reorganised, while the body
imposes a kinematic ``blocking'' that suppresses the wall-normal velocity
component close to the surface even as the straining amplifies it farther
upstream \citep{HuntGraham1978,GoldsteinAtassi1976}. The upwash spectrum
\emph{seen by the leading edge} is therefore distorted, scale dependent
and tensorially redistributed relative to its upstream counterpart. This
pre-impact distortion stage lies precisely between the upstream
measurement plane and the input required by Amiet's theory, and it is the
stage that the standard prediction chain omits.

Several routes to recovering this distortion exist, each with a
characteristic limitation. Scale-resolving simulation---large-eddy
simulation, lattice-Boltzmann methods or direct computational
aeroacoustics---represents the distortion and the radiation directly, but
at a cost that precludes broad parametric study across turbulence
intensity, length-scale ratio and Reynolds number. Linear rapid-distortion
theory (RDT) provides an analytical description of the mean-strain
distortion and, in its inhomogeneous form, of the surface blocking
\citep{HuntCarruthers1990,GoldsteinAfsarLeib2013}; it has been used to
construct RDT-modified turbulence spectra that are fed to Amiet's solution
in place of the isotropic input \citep{deSantana2016}. RDT is, however,
linear, and in its tractable homogeneous form it neglects the nonlinear
inter-scale transfer and the slow return-to-isotropy that act over the
finite approach time. Synthetic-turbulence methods---synthetic-eddy and
random-particle techniques---efficiently reproduce a \emph{prescribed}
spectrum and coherence and are now routine for generating aeroacoustic
inflow \citep{KimHaeri2015}, but by construction they impose the spectrum
rather than evolving it, and so cannot themselves predict how the spectrum
is modified before impact. What is missing is a reduced-order model that
\emph{evolves} the incoming turbulence through the distortion region with
full scale resolution and at negligible cost, delivering the distorted
upwash spectrum as an output rather than an assumption.

The one-dimensional turbulence (ODT) model of \citet{Kerstein1999} is a
natural candidate for this role. ODT represents the evolution of a
turbulent velocity profile along a single notional line of sight,
resolving the full range of length scales in one spatial dimension while
modelling turbulent advection as a stochastic sequence of scale-local
mapping (``eddy'') events built on the triplet map introduced in the
linear-eddy model \citep{Kerstein1991}. Its vector formulation
\citep{Kerstein2001} carries all three velocity components and
redistributes energy among them through a kernel mechanism that emulates
pressure scrambling \citep{WunschKerstein2001,Kerstein2001}, reproducing
inter-component transfer and return to isotropy; the model has been generalised to variable density
\citep{AshurstKerstein2005}, to efficient adaptive-mesh and spatially
developing implementations \citep{Lignell2013,StephensLignell2021}, to
cylindrical and spherical geometries \citep{Lignell2018}, and to a
three-dimensional large-eddy framework \citep{SchmidtKerstein2010}. ODT
reproduces inertial-range scaling and component statistics of free-shear
and wall-bounded flows at a small fraction of the cost of
three-dimensional simulation. Of particular relevance here, ODT velocity
fields have recently been coupled to acoustic analogies to estimate the
far-field sound of low-Mach-number turbulent jets
\citep{SharmaKleinSchmidt2022,MedinaMendez2023}, demonstrating that ODT
statistics---including the component spectra and fluctuating Reynolds
stresses---can be supplied to an external acoustic theory. These
developments establish ODT as an accurate and inexpensive generator of
scale-resolved turbulence statistics, and in particular of the
wall-normal velocity spectrum required by Amiet's theory.

Standard ODT cannot, however, represent the pre-impact distortion itself.
The model contains no mechanism for the action of an imposed mean strain
on the fluctuating field: its pressure treatment, embodied in the
scrambling kernel, models only the slow, turbulence-driven return to
isotropy and not the rapid, mean-strain-driven redistribution that governs
the stagnation region. The apparent remedy---importing the linear RDT
amplitude equation, in which the pressure response is the solenoidal
projection $P_{ij}(\boldsymbol{k}) = \delta_{ij} - k_i k_j/k^2$ acting in
three-dimensional wavenumber space---is fundamentally incompatible with
the one-dimensional model, because the projection requires the full
wavevector $\boldsymbol{k}=(k_1,k_2,k_3)$, whereas a single ODT line
resolves only the wavenumber along its own coordinate. Reconciling the
inherently three-dimensional and non-local pressure physics of rapid
distortion with the one-dimensional and local structure of ODT is the
central modelling problem addressed in this paper.

We present a strain-coupled formulation of one-dimensional turbulence that
incorporates rapid distortion by an imposed mean velocity gradient while
preserving the structural foundations of the model. The formulation rests
on a decomposition of the pressure--strain interaction into its rapid and
slow parts \citep{Pope2000}: the slow, turbulence-driven part is retained
as the existing ODT scrambling kernel, while the rapid, mean-strain-driven
part is introduced as a new component-redistribution kernel. Rather than
evaluating the three-dimensional projection, the rapid kernel is
constructed so that the component-energy evolution of the line reproduces
the exact homogeneous-RDT solution for an imposed mean strain, rendering
the formulation provably RDT-consistent in the homogeneous limit while
remaining expressed entirely in quantities available on the line. The
mean-strain production enters as a continuous forcing of the velocity
components, which we show preserves the measure-preservation, continuity
and conservation properties of the triplet-map mechanism; the kinematic
compression of wavenumbers under strain is captured by a consistent
deformation of the ODT domain, connecting the formulation to its
adaptive-mesh implementation. The result is a reduced-order, scale-resolved
model of the pre-impact distortion stage whose output---the distorted,
anisotropic upwash spectrum---is intended to replace the prescribed
isotropic input of Amiet-type prediction, with the airfoil response, edge
scattering and spanwise coherence retained as the classical, downstream
acoustic stage.

The remainder of the paper is organised as follows.
Section~\ref{sec:odt} reviews the elements of standard ODT required for the
present development and the rapid/slow structure of the pressure--strain
interaction. Section~\ref{sec:formulation} derives the strain-coupled
formulation, establishes the preservation of the model's structural
invariants under mean-strain forcing, and constructs the RDT-consistent
rapid kernel. Section~\ref{sec:validation} validates the formulation
against analytical homogeneous-RDT solutions for the canonical strains
relevant to the stagnation region---axisymmetric contraction and plane
strain---and against reference scale-resolving data. Section~\ref{sec:upwash}
demonstrates the resulting distorted upwash spectrum and its implications
for the Amiet input. Section~\ref{sec:conclusions} concludes and outlines
the remaining steps---surface blocking and spanwise coherence---towards a
complete leading-edge-noise framework.

% =====================================================================
%  Section 2  --  Background
%  (i) standard one-dimensional turbulence; (ii) rapid/slow structure
%      of the pressure--strain interaction and linear RDT.
%  Target: Journal of Fluid Mechanics (jfm.cls); citations in references.bib
% =====================================================================

\section{Background}\label{sec:odt}

This section collects the elements on which the strain-coupled formulation
of \S\,\ref{sec:formulation} is built. Section~\ref{sec:odt-standard}
summarises standard one-dimensional turbulence (ODT) in the notation used
throughout; \S\,\ref{sec:rapid-slow} recalls the rapid/slow decomposition
of the pressure--strain interaction and the linear rapid-distortion
equations that the new formulation is required to reproduce. Nothing in
this section is original; readers familiar with
\citet{Kerstein1999,Kerstein2001} and with second-moment closure
\citep{Pope2000} may proceed to \S\,\ref{sec:formulation} after noting the
notation and the correspondence stated in \S\,\ref{sec:rapid-slow}.

\subsection{Standard one-dimensional turbulence}\label{sec:odt-standard}

ODT evolves a three-component velocity field $u_i(y,t)$, $i=1,2,3$, and any
number of scalar fields $\theta(y,t)$, defined on a one-dimensional domain
parameterised by the spatial coordinate $y$, which is identified with the
coordinate direction $x_2$ \citep{Kerstein2001}. The fields represent
individual realisations: ensemble statistics are accumulated over many
independent simulations rather than carried by the fields themselves. The
fields evolve by two interleaved mechanisms---continuous molecular
transport and a stochastic sequence of instantaneous \emph{eddy events}.

\subsubsection{Molecular transport}

Between events the fields obey ordinary one-dimensional diffusion,
\begin{equation}
  \left(\partial_t-\nu\,\partial_y^2\right)u_i(y,t)=0,
  \qquad
  \left(\partial_t-\kappa\,\partial_y^2\right)\theta(y,t)=0,
  \label{eq:diffusion}
\end{equation}
with kinematic viscosity $\nu$ and scalar diffusivity $\kappa$
(Schmidt number $\mathit{Sc}=\nu/\kappa$). In a constant-property,
constant-density medium this is the only continuous evolution; the
variable-density generalisation is given by \citet{AshurstKerstein2005}.

\subsubsection{Eddy events: the triplet map and kernel}

An eddy event models the effect of a single turbulent overturn. Because a
continuum rotational motion cannot be represented on a line, the event is
implemented as an instantaneous, measure-preserving rearrangement of the
profiles together with an additive velocity change,
\begin{equation}
  u_i(y)\;\longrightarrow\;u_i\!\big(f(y)\big)+c_i\,K(y),
  \qquad
  \theta(y)\;\longrightarrow\;\theta\!\big(f(y)\big).
  \label{eq:event}
\end{equation}
The map $f$ is the \emph{triplet map} \citep{Kerstein1991,Kerstein1999},
defined on an interval $[y_0,y_0+l]$ of size $l$ by
\begin{equation}
  f(y)=y_0+
  \begin{cases}
    3(y-y_0), & y_0\le y\le y_0+\tfrac{1}{3}l,\\[2pt]
    2l-3(y-y_0), & y_0+\tfrac{1}{3}l\le y\le y_0+\tfrac{2}{3}l,\\[2pt]
    3(y-y_0)-2l, & y_0+\tfrac{2}{3}l\le y\le y_0+l,\\[2pt]
    y-y_0, & \text{otherwise,}
  \end{cases}
  \label{eq:triplet}
\end{equation}
written here through its inverse, so that fluid at $f(y)$ is placed at $y$.
The map compresses the interval to one-third of its length, inserts three
copies, and reverses the central copy; it leaves the field outside
$[y_0,y_0+l]$ unchanged. The triplet map satisfies the three structural
requirements that underpin the model: it is \emph{measure preserving} (the
non-local analogue of a divergence-free rearrangement), \emph{continuous}
(it introduces no discontinuities into the fields), and \emph{scale local}
(it changes property gradients by at most a factor of order unity)
\citep{Kerstein2001}. These three properties are exactly the invariants
whose preservation under an imposed mean strain is established in
\S\,\ref{sec:formulation}.

The additive term in \eqref{eq:event} is built from the \emph{kernel}
\begin{equation}
  K(y)\equiv y-f(y),
  \label{eq:kernel}
\end{equation}
which is non-zero only within the eddy interval and satisfies the two
identities used below,
\begin{equation}
  \int K(y)\,\mathrm{d}y = 0,
  \qquad
  \int K^2(y)\,\mathrm{d}y = \tfrac{4}{27}\,l^{3}.
  \label{eq:kernel-identities}
\end{equation}
The first guarantees that the additive change conserves the $y$-integrated
momentum of each velocity component at constant density; the second fixes
the energetics of the redistribution. The amplitudes $c_i$ are determined
event-by-event by the pressure-scrambling rule described next.

\subsubsection{Pressure scrambling and component-energy redistribution}

The additive kernel term in \eqref{eq:event} is the model surrogate for the
pressure-mediated redistribution of energy among velocity components, an
effect absent from the original single-component formulation
\citep{Kerstein1999} and introduced through the kernel by
\citet{WunschKerstein2001,Kerstein2001}. Using measure preservation and
\eqref{eq:kernel-identities}, the change in the kinetic energy of component
$i$ produced by an event is
\begin{equation}
  \Delta E_i = \rho_0\, l^2\, c_i\!\left(u_{i,K}+\tfrac{2}{27}\,l\,c_i\right),
  \qquad
  u_{i,K}\equiv\frac{1}{l^2}\int u_i\!\big(f(y)\big)\,K(y)\,\mathrm{d}y,
  \label{eq:dEi}
\end{equation}
where $\rho_0$ is the (constant) density per unit length, which does not
affect the constant-density applications considered here. The amplitudes
are fixed by requiring that the event redistribute energy among components
while conserving the total, $\sum_i\Delta E_i=0$, and that the rule be
invariant under permutation of the component indices. These requirements
give
\begin{equation}
  \Delta E_i = \alpha\sum_j T_{ij}\,Q_j,
  \qquad
  \mathsf{T}=\frac{1}{2}
  \begin{pmatrix}
    -2 & 1 & 1\\ 1 & -2 & 1\\ 1 & 1 & -2
  \end{pmatrix},
  \qquad
  Q_i\equiv\tfrac{27}{8}\,\rho_0\, l\, u_{i,K}^2,
  \label{eq:redistribution}
\end{equation}
in which $Q_i$ is the \emph{available kinetic energy} of component $i$
within the eddy and $\alpha\in[0,1]$ is the model parameter setting the
fraction of extractable energy transferred towards isotropy
\citep{Kerstein2001}. Combining \eqref{eq:dEi} and
\eqref{eq:redistribution} determines the amplitudes,
\begin{equation}
  c_i = \frac{27}{4l}\left[-u_{i,K}
        +\operatorname{sgn}(u_{i,K})
        \sqrt{u_{i,K}^2+\alpha\sum_j T_{ij}\,u_{j,K}^2}\,\right],
  \label{eq:ci}
\end{equation}
which is real for $0\le\alpha\le1$ and vanishes as $\alpha\to0$. The
transfer matrix $\mathsf{T}$ acts to equalise the component energies; the
construction \eqref{eq:redistribution}--\eqref{eq:ci} is therefore a model
of the \emph{return-to-isotropy} (slow) part of the pressure--strain
interaction, a point made precise in \S\,\ref{sec:rapid-slow}. Standard ODT
contains no counterpart of the rapid, mean-strain-driven part.

\subsubsection{Eddy selection}

The event sequence is a Poisson process whose rate reflects the
instantaneous state of the flow. Each candidate eddy $(y_0,l)$ is assigned
a time scale $\tau(y_0,l;t)$ built from the available energy of the
velocity component aligned with the domain,
\begin{equation}
  \left(\frac{l}{\tau}\right)^{2}
  \sim\; u_{2,K}^2+\alpha\sum_j T_{2j}\,u_{j,K}^2-Z\,\frac{\nu^2}{l^2},
  \label{eq:timescale}
\end{equation}
the final term being a viscous penalty, controlled by the order-unity
parameter $Z$, that suppresses eddies slower than their viscous time. The
corresponding event-rate distribution is
\begin{equation}
  \lambda(y_0,l;t)=\frac{C}{l^2\,\tau(y_0,l;t)},
  \label{eq:rate}
\end{equation}
with rate constant $C$; eddies for which the right-hand side of
\eqref{eq:timescale} is negative are forbidden. Crucially, the
inertial-range cascade and the $E(k)\sim k^{-5/3}$ scaling are
\emph{outcomes} of repeated sampling from \eqref{eq:rate}, not inputs to
it \citep{Kerstein1999}. The three constants $C$, $Z$ and a large-eddy
suppression parameter are the only tunable inputs of the model
\citep{StephensLignell2021}.

\subsubsection{Statistical observables}\label{sec:observables}

ODT yields single- and multi-point moments, probability density functions,
conditional statistics, power spectra and Lagrangian statistics
\citep{Kerstein1999}. For the present work the relevant observable is the
single-component velocity spectrum---in particular the spectrum and
variance of the component normal to the airfoil, which are second moments
of the velocity field and are obtained directly from the realisations. A
distinction emphasised by \citet{Kerstein2001} must, however, be respected:
because the line velocity does not itself advect fluid, an off-diagonal
Reynolds stress such as $\langle u_1' u_2'\rangle$ \emph{cannot} be
interpreted as a momentum flux when formed as a product of the line
velocity components; such advective fluxes are instead evaluated by
monitoring the transport induced by the eddy events. This distinction
constrains how the anisotropy produced by the formulation of
\S\,\ref{sec:formulation} is to be diagnosed and validated
(\S\,\ref{sec:validation}): diagonal component spectra are read from the
fields, whereas shear stresses are read from eddy-induced fluxes.

\subsection{Rapid and slow pressure--strain, and linear rapid
distortion}\label{sec:rapid-slow}

The Reynolds-stress transport equation for a fluctuating field $u_i'$ about
a mean $U_i$ with constant velocity gradient $A_{ij}\equiv\partial
U_i/\partial x_j$ reads, in homogeneous turbulence,
\begin{equation}
  \frac{\mathrm{d}\langle u_i' u_j'\rangle}{\mathrm{d}t}
  = \mathcal{P}_{ij} + \Pi_{ij} - \varepsilon_{ij},
  \label{eq:rstransport}
\end{equation}
with production, pressure--strain and dissipation
\begin{equation}
  \mathcal{P}_{ij}=-\langle u_i' u_k'\rangle A_{jk}
                   -\langle u_j' u_k'\rangle A_{ik},
  \quad
  \Pi_{ij}=\Big\langle \tfrac{p'}{\rho}
            \big(\partial_j u_i'+\partial_i u_j'\big)\Big\rangle,
  \quad
  \varepsilon_{ij}=2\nu\langle \partial_k u_i'\,\partial_k u_j'\rangle .
  \label{eq:Pij}
\end{equation}
Taking the divergence of the fluctuating momentum equation gives the
pressure as the solution of a Poisson equation whose source splits into a
part linear in the mean velocity gradient and a part quadratic in the
fluctuations,
\begin{equation}
  \frac{1}{\rho}\nabla^2 p'
  = \underbrace{-2\,A_{lm}\,\partial_l u_m'}_{\text{rapid}}
    \;-\;
    \underbrace{\partial_l\partial_m\!\left(u_l' u_m'
     -\langle u_l' u_m'\rangle\right)}_{\text{slow}} .
  \label{eq:poisson}
\end{equation}
The pressure--strain correlation inherits this decomposition,
$\Pi_{ij}=\Pi_{ij}^{(r)}+\Pi_{ij}^{(s)}$, into a \emph{rapid} part driven
by the mean strain and a \emph{slow} part driven by the turbulence alone
\citep{Pope2000}. The slow part relaxes the Reynolds stress towards
isotropy; the simplest closure is that of Rotta,
$\Pi_{ij}^{(s)}=-C_1(\varepsilon/k)\big(\langle u_i'u_j'\rangle
-\tfrac{2}{3}k\,\delta_{ij}\big)$, with $k=\tfrac12\langle u_i'u_i'\rangle$.
It is this return-to-isotropy mechanism that the ODT kernel
\eqref{eq:redistribution}--\eqref{eq:ci} reproduces in model form. The
rapid part is linear in $A_{lm}$ and is determined, in homogeneous
turbulence, by the velocity-spectrum tensor
$\Phi_{ij}(\boldsymbol{k})$; the corresponding ``isotropisation of
production'' closure is
$\Pi_{ij}^{(r)}=-C_2\big(\mathcal{P}_{ij}-\tfrac{2}{3}\mathcal{P}\,
\delta_{ij}\big)$, with $\mathcal{P}=\tfrac12\mathcal{P}_{ii}$
\citep{Pope2000}. Standard ODT represents neither $\mathcal{P}_{ij}$ nor
$\Pi_{ij}^{(r)}$.

For a rapidly applied strain---formally, when the total strain accumulates
in a time short compared with the eddy-turnover time---the rapid part can
be computed exactly from linear theory \citep{BatchelorProudman1954,
Hunt1973,HuntCarruthers1990}. Writing the fluctuation as a sum of Fourier
modes that follow the mean flow, each mode of (time-dependent) wavevector
$\boldsymbol{k}(t)$ and amplitude $\hat{u}_i(\boldsymbol{k},t)$ evolves,
in the inviscid rapid limit, according to
\begin{equation}
  \frac{\mathrm{d}k_i}{\mathrm{d}t}=-A_{ji}\,k_j,
  \qquad
  \frac{\mathrm{d}\hat{u}_i}{\mathrm{d}t}
  =\left(\frac{2\,k_i k_m}{k^2}-\delta_{im}\right)A_{mn}\,\hat{u}_n ,
  \label{eq:rdt}
\end{equation}
where $k^2=k_m k_m$. The bracketed operator combines the production
$-A_{im}\hat{u}_m$ with the rapid pressure response, the latter being the
solenoidal projection
$P_{ij}(\boldsymbol{k})=\delta_{ij}-k_ik_j/k^2$ applied to the strained
field. The factor of two in \eqref{eq:rdt} is fixed by the requirement that
the motion remain divergence free: using both relations in
\eqref{eq:rdt},
\begin{equation}
  \frac{\mathrm{d}}{\mathrm{d}t}\big(k_i\hat{u}_i\big)
  = \dot{k}_i\hat{u}_i + k_i\dot{\hat{u}}_i
  = -A_{ji}k_j\hat{u}_i
    + \Big(\tfrac{2k^2 k_m}{k^2}-k_m\Big)A_{mn}\hat{u}_n
  = 0,
  \label{eq:solenoidal}
\end{equation}
so that $k_i\hat{u}_i=0$ is preserved; any other coefficient violates
continuity. Equations~\eqref{eq:rdt} integrate the spectral tensor through
the distortion and thereby determine the rapid pressure--strain
$\Pi_{ij}^{(r)}$ and the component-energy evolution exactly.

Two features of \eqref{eq:rdt} are central to what follows. First, the
pressure response is the operator $P_{ij}(\boldsymbol{k})$, which is
\emph{non-local} (it couples the components through the full wavevector)
and intrinsically \emph{three-dimensional} (it requires all of $k_1,k_2,k_3$).
A single ODT line resolves only the wavenumber along its own coordinate,
so \eqref{eq:rdt} cannot be evaluated on the line as written; reconciling
this incompatibility is the subject of \S\,\ref{sec:formulation}. Second,
the decomposition $\Pi_{ij}=\Pi_{ij}^{(r)}+\Pi_{ij}^{(s)}$ furnishes the
organising principle of the new formulation: the slow part is already
carried by the ODT kernel \eqref{eq:redistribution}--\eqref{eq:ci}, and it
remains only to supply the production $\mathcal{P}_{ij}$ and a rapid
redistribution consistent with \eqref{eq:rdt}, expressed in quantities
available on the line.

% =====================================================================
%  Section 3  --  Strain-coupled one-dimensional turbulence
%  The new formulation: production forcing, RDT-consistent rapid kernel,
%  domain straining. Self-contained derivations.
%
%  Target: Journal of Fluid Mechanics (jfm.cls).
%  Requires in preamble:
%     \usepackage{amsmath,amssymb,amsthm}
\newtheorem{proposition}{Proposition}[section]
%  Citations resolve against references.bib.
% =====================================================================

\section{Strain-coupled one-dimensional turbulence}\label{sec:formulation}

The linear rapid-distortion equations~\eqref{eq:rdt} contain two distinct
ingredients: an evolution of the Fourier \emph{amplitudes}
$\hat{u}_i(\boldsymbol{k},t)$, comprising production and the solenoidal
pressure projection, and an evolution of the \emph{wavevectors}
$\boldsymbol{k}(t)$ that carries the kinematic straining of scales. The
amplitude equation is the obstruction identified in
\S\,\ref{sec:rapid-slow}: its pressure term requires the full
three-dimensional wavevector and is therefore not directly representable on
a one-dimensional line. Our strategy is to split the distortion physics
according to this structure and to assign each part to the ODT mechanism
that can carry it faithfully. The production, which is local in physical
space, is imposed as a continuous forcing of the line velocity
(\S\,\ref{sec:strain-forcing}); the rapid pressure--strain, which is not,
is supplied as a continuous, energy-conserving redistribution operator
constructed to reproduce the homogeneous rapid-distortion evolution of the
component energies under an explicit spectral closure
(\S\,\ref{sec:rapid-kernel}); and the wavevector kinematics are recovered by
a consistent straining of the ODT domain
(\S\,\ref{sec:domain-strain}). The slow, turbulence-driven
pressure--strain and the inertial cascade remain the discrete eddy events
of standard ODT, unchanged. Section~\ref{sec:summary} assembles the
formulation and identifies the parameter governing the balance between the
rapid and slow mechanisms.

Throughout, $A_{ij}=\partial U_i/\partial x_j$ denotes the imposed mean
velocity gradient, taken uniform across the line at each instant and
varying along the mean trajectory; the medium is incompressible,
$A_{ii}=0$. The line carries the fluctuation field $u_i(y,t)$ about this
mean. We write $R_{ij}\equiv\overline{u_i u_j}$ for the Reynolds-stress
tensor, where the overline denotes the average over the statistically
homogeneous directions (equivalently, the ensemble of independent
realisations), and $k_t\equiv\tfrac12 R_{ii}$ for the turbulent kinetic
energy. In the homogeneous configurations used to establish and test the
formulation (\S\,\ref{sec:validation}) the spatial line average and the
ensemble average coincide.

\subsection{The strain-coupled line equation and its structural
admissibility}\label{sec:strain-forcing}

In rapid-distortion theory the production term of the fluctuation momentum
equation, $-A_{ij}u_j$, is a pointwise linear combination of the local
velocity components and is therefore exactly representable on the line.
Between eddy events we replace the molecular evolution
\eqref{eq:diffusion} by
\begin{equation}
  \partial_t u_i
  = \nu\,\partial_y^2 u_i \;-\; A_{ij}\,u_j \;+\; B_{ij}\,u_j,
  \label{eq:cont-evolution}
\end{equation}
in which $-A_{ij}u_j$ is the production and $B_{ij}u_j$ is the rapid
pressure--strain operator constructed in \S\,\ref{sec:rapid-kernel}; we
defer $B_{ij}$ and treat the production term first. The eddy events
\eqref{eq:event}--\eqref{eq:ci} are retained without modification. The
evolution is thus advanced by operator splitting between a continuous phase
governed by \eqref{eq:cont-evolution} and a discrete sequence of events
sampled from \eqref{eq:rate}.

A formulation is admissible only if the additional forcing does not
violate the structural properties on which ODT rests. We establish this
directly.

\begin{proposition}[Structural admissibility]\label{prop:admissibility}
The strain forcing in \eqref{eq:cont-evolution} preserves the four
structural properties of standard ODT: \emph{(i)} the measure-preservation,
\emph{(ii)} continuity and \emph{(iii)} scale-locality of the triplet map,
and \emph{(iv)} the kernel identities~\eqref{eq:kernel-identities}.
\end{proposition}

\begin{proof}
Properties (i)--(iv) are properties of the discrete event, that is, of the
map $f$~\eqref{eq:triplet} and the kernel $K$~\eqref{eq:kernel}. Because the
advancement is split into a continuous phase \eqref{eq:cont-evolution} and
the unaltered discrete event, the map and kernel are applied exactly as in
standard ODT; properties (i), (iii) and (iv), which are geometric
statements about $f$ and $K$ on a fixed domain, therefore hold verbatim.
It remains to verify that the continuous phase does not destroy the
continuity of the fields, which underlies (ii). For a mean gradient uniform
along the line, the linear, $y$-independent operator
$\mathcal{A}_{ij}\equiv -A_{ij}+B_{ij}$ acts pointwise, and the
solution of the advective part of \eqref{eq:cont-evolution} over a
sub-step $[t_0,t]$ is
\begin{equation}
  u_i(y,t)=\big[\exp\!\textstyle\int_{t_0}^{t}\mathcal{A}\,\mathrm{d}t'\big]_{ij}\,
           u_j(y,t_0),
  \label{eq:matrixexp}
\end{equation}
a matrix that is independent of $y$. It maps continuous (respectively,
$C^\infty$) profiles to continuous ($C^\infty$) profiles and introduces no
spatial discontinuity, while the molecular term $\nu\partial_y^2 u_i$ is
smoothing. Continuity is preserved, establishing (ii). The forcing
displaces no fluid along $y$, so it leaves the Lebesgue measure of the
domain unchanged; the deliberate straining of the domain is treated
separately in \S\,\ref{sec:domain-strain}.
\end{proof}

The energetic role of the forcing follows from \eqref{eq:cont-evolution}.
At fixed domain length, the production term evolves the Reynolds stress as
\begin{equation}
  \left.\frac{\mathrm{d}R_{ij}}{\mathrm{d}t}\right|_{\text{prod}}
  = \overline{(\dot u_i u_j + u_i \dot u_j)}
  = -A_{ik}R_{kj}-A_{jk}R_{ki}
  \;\equiv\; \mathcal{P}_{ij},
  \label{eq:production-rate}
\end{equation}
which is \emph{exactly} the Reynolds-stress production tensor of
\eqref{eq:Pij}; the corresponding kinetic-energy production is
$\mathcal{P}\equiv\tfrac12\mathcal{P}_{ii}=-A_{ij}R_{ij}$. The production
term thus carries no model assumption. With diffusion supplying
dissipation through \eqref{eq:diffusion}, the kinetic-energy budget of the
formulation separates cleanly into an exact production by the mean strain,
a redistribution that conserves energy (the kernels), and a dissipation
(molecular diffusion). This is the structure of the exact Reynolds-stress
equation \eqref{eq:rstransport}, and it is the reason the forcing is
admissible: it adds the one term---production---that standard ODT lacks,
without disturbing the terms it already represents.

\subsection{The RDT-consistent rapid redistribution
operator}\label{sec:rapid-kernel}

\subsubsection{Why production alone is insufficient}

Were the continuous evolution to consist of production and diffusion only
($B_{ij}=0$), the line would reproduce $\mathrm{d}R_{ij}/\mathrm{d}t=
\mathcal{P}_{ij}$, whereas the exact inviscid rapid-distortion balance is
$\mathrm{d}R_{ij}/\mathrm{d}t=\mathcal{P}_{ij}+\Pi_{ij}^{(r)}$. The
omission is not small. For turbulence that is isotropic at the onset of
straining, the classical result derived below gives
$\Pi_{ij}^{(r)}=-\tfrac35\mathcal{P}_{ij}$, so that the true initial
distortion rate is $\tfrac25\mathcal{P}_{ij}$: production acting alone
overstates the rate at which anisotropy develops by a factor of
$\tfrac52$. For sustained strong strain, neglect of the pressure response
also permits a component energy to grow without the bound that
incompressibility imposes, violating realisability. Representing
$\Pi_{ij}^{(r)}$ is therefore essential, not optional.

\subsubsection{The exact rapid pressure--strain and the closure problem}

The exact rapid pressure--strain follows from the amplitude
equation~\eqref{eq:rdt}. Writing the operator there as
$\mathrm{d}\hat{u}_i/\mathrm{d}t = M_{in}\hat{u}_n$ with
$M_{in}=-A_{in}+2(k_ik_m/k^2)A_{mn}$, and noting that the incompressible
wavevector dynamics $\dot{k}_i=-A_{ji}k_j$ preserve the measure
$\mathrm{d}\boldsymbol{k}$ (its phase-space divergence is $-A_{ii}=0$), the
spectral tensor $\Phi_{ij}=\langle \hat{u}_i^{*}\hat{u}_j\rangle$ yields
\begin{equation}
  \frac{\mathrm{d}R_{ij}}{\mathrm{d}t}
  = \int\!\big(M_{in}\Phi_{nj}+M_{jn}\Phi_{in}\big)\,
    \mathrm{d}\boldsymbol{k}
  = \mathcal{P}_{ij} + \Pi_{ij}^{(r)},
  \label{eq:dRdt-rdt}
\end{equation}
in which the production~\eqref{eq:production-rate} has been separated from
the pressure contribution,
\begin{equation}
  \Pi_{ij}^{(r)} = 2A_{mn}\big(\mathcal{R}_{ijmn}+\mathcal{R}_{jimn}\big),
  \qquad
  \mathcal{R}_{ijmn} \equiv
  \int \frac{k_i k_m}{k^2}\,\Phi_{nj}(\boldsymbol{k})\,
  \mathrm{d}\boldsymbol{k}.
  \label{eq:exact-rps}
\end{equation}
The tensor $\Pi_{ij}^{(r)}$ is traceless: using the solenoidality
$k_n\Phi_{nj}=0$ (a consequence of $k_i\hat{u}_i=0$),
\begin{equation}
  \Pi_{ii}^{(r)} = 4A_{mn}\!\int\frac{k_i k_m}{k^2}\Phi_{ni}\,
  \mathrm{d}\boldsymbol{k}
  = 4A_{mn}\!\int\frac{k_m}{k^2}\big(k_i\Phi_{ni}\big)\,
  \mathrm{d}\boldsymbol{k}=0,
  \label{eq:traceless}
\end{equation}
so the rapid pressure--strain redistributes energy among components
without changing the total, in keeping with the energy-conserving
character of the ODT kernel mechanism.

Equation~\eqref{eq:exact-rps} exposes the central difficulty. The rapid
pressure--strain is a functional of the spectral tensor $\Phi_{nj}
(\boldsymbol{k})$ through the orientation-weighted integral
$\mathcal{R}_{ijmn}$; it is \emph{not} expressible in terms of the
single-point Reynolds stress $R_{ij}$ alone. This is the well-known closure
problem of the rapid pressure--strain, common to every single-point model
\citep{Crow1968,LaunderReeceRodi1975,Pope2000}; it is intrinsic to the
physics and not a deficiency peculiar to the present approach. A
one-dimensional line carries $R_{ij}$ and the one-dimensional spectrum
$\Phi_{ij}(k_2)$ along its own coordinate, but it does not carry the
distribution of energy over wavevector \emph{orientation} that
$\mathcal{R}_{ijmn}$ requires. A spectral-tensor closure is therefore
unavoidable, exactly as in second-moment closure; we make it explicit.

\subsubsection{Baseline isotropic closure}

The minimal closure assumes the spectral tensor to be instantaneously
isotropic,
$\Phi_{nj}(\boldsymbol{k})=\frac{E(k)}{4\pi k^2}\big(\delta_{nj}
-k_nk_j/k^2\big)$ with $\int E(k)\,\mathrm{d}k=k_t$. Writing
$e_i\equiv k_i/k$ and using the angular averages over the unit sphere
$\overline{e_i e_m}=\tfrac13\delta_{im}$ and
$\overline{e_i e_m e_n e_j}=\tfrac1{15}(\delta_{im}\delta_{nj}
+\delta_{in}\delta_{mj}+\delta_{ij}\delta_{mn})$,
\begin{equation}
  \mathcal{R}_{ijmn}^{\,\mathrm{iso}}
  = k_t\Big[\tfrac{4}{15}\delta_{im}\delta_{nj}
    -\tfrac1{15}\big(\delta_{in}\delta_{mj}+\delta_{ij}\delta_{mn}\big)\Big].
  \label{eq:iso-closure}
\end{equation}
Substituting \eqref{eq:iso-closure} into \eqref{eq:exact-rps} and using
$A_{nn}=0$ gives, with $S_{ij}\equiv\tfrac12(A_{ij}+A_{ji})$,
\begin{equation}
  \Pi_{ij}^{(r),\mathrm{iso}}
  = 2k_t\Big[\tfrac{4}{15}A_{ij}-\tfrac1{15}A_{ji}\Big]
   +2k_t\Big[\tfrac{4}{15}A_{ji}-\tfrac1{15}A_{ij}\Big]
  = \tfrac{4}{5}\,k_t\,S_{ij}.
  \label{eq:Pir-iso}
\end{equation}
For isotropic $R_{ij}=\tfrac23 k_t\delta_{ij}$ the production
\eqref{eq:production-rate} is $\mathcal{P}_{ij}=-\tfrac43 k_t S_{ij}$, so
that \eqref{eq:Pir-iso} is equivalent to
\begin{equation}
  \Pi_{ij}^{(r),\mathrm{iso}} = -\tfrac{3}{5}\,\mathcal{P}_{ij},
  \label{eq:crow}
\end{equation}
the exact response of initially isotropic turbulence to a weak rapid strain
\citep{Crow1968}. Equation~\eqref{eq:crow} is the content of the
isotropisation-of-production model with the constant $C_2=\tfrac35$
\citep{Pope2000}; here it is obtained, with no adjustable constant, as the
isotropic limit of the rapid-distortion integral~\eqref{eq:exact-rps}.

When the turbulence is anisotropic, \eqref{eq:iso-closure} no longer holds
and $\mathcal{R}_{ijmn}$ must be modelled as a tensor function of $R_{ij}$;
the general representation linear in the anisotropy is that of
\citet{LaunderReeceRodi1975}, which introduces calibrated constants and to
which the present construction reduces in form. We adopt
\eqref{eq:iso-closure} as the baseline and treat the anisotropic
correction as a modelling option, returning to its assessment in
\S\,\ref{sec:validation}.

\subsubsection{Realisation on the line}

Given the closed rapid pressure--strain $\Pi_{ij}^{(r)}$---symmetric and
traceless by \eqref{eq:traceless}---we require a pointwise operator
$B_{ij}$ on the line whose action reproduces it. Demanding that the
operator evolve the Reynolds stress by exactly $\Pi_{ij}^{(r)}$,
\begin{equation}
  B_{ik}R_{kj}+B_{jk}R_{ki}
  = \big(\mathsf{B}\mathsf{R}+\mathsf{R}\mathsf{B}\big)_{ij}
  = \Pi_{ij}^{(r)},
  \label{eq:lyapunov}
\end{equation}
where the second equality holds for symmetric $\mathsf{B}$. Equation
\eqref{eq:lyapunov} is a continuous Lyapunov equation; for a
positive-definite Reynolds stress $\mathsf{R}\succ0$ it possesses a unique
symmetric solution $\mathsf{B}$. The solution conserves kinetic energy
automatically: contracting \eqref{eq:lyapunov} gives
$2B_{ij}R_{ij}=\Pi_{ii}^{(r)}=0$, so that the rate of working of the
operator on the line, $\overline{B_{ij}u_iu_j}=B_{ij}R_{ij}$, vanishes. In
the isotropic case $\mathsf{R}=\tfrac23 k_t\mathsf{I}$, \eqref{eq:lyapunov}
together with \eqref{eq:Pir-iso} reduces to $\tfrac43 k_t\mathsf{B}
=\tfrac45 k_t\mathsf{S}$, that is, the closed form
\begin{equation}
  B_{ij} = \tfrac{3}{5}\,S_{ij},
  \qquad
  \mathcal{A}_{ij} = -A_{ij} + \tfrac{3}{5}\,S_{ij}.
  \label{eq:iso-B}
\end{equation}
The operator $B_{ij}$ is evaluated from the running Reynolds stress; it is
applied continuously, paced by the mean strain, exactly as is the
production. This is deliberate: the rapid pressure--strain is, by
construction, the part of the redistribution driven by the mean velocity
gradient and therefore acts at the strain rate, in contrast to the slow,
turbulence-driven redistribution that the eddy kernel
\eqref{eq:redistribution}--\eqref{eq:ci} supplies at the eddy rate. The two
mechanisms are thereby cleanly separated: strain-driven effects (production
and rapid pressure--strain) are continuous and enter through the operator
$\mathcal{A}_{ij}$, while turbulence-driven effects (the inertial cascade
and the slow return to isotropy) remain the discrete events. The uniform,
scale-independent action of $B_{ij}$ is the faithful single-point reduction
of the rapid pressure response, which in \eqref{eq:rdt} acts on every
Fourier mode through the projection $P_{ij}(\boldsymbol{k})$; the
distribution of energy across scales continues to be governed by diffusion,
the triplet-map cascade and the domain straining of
\S\,\ref{sec:domain-strain}.

We can now state precisely the sense in which the formulation is
rapid-distortion consistent.

\begin{proposition}[Rapid-distortion consistency]\label{prop:rdt}
Let $B_{ij}$ be defined by the Lyapunov equation~\eqref{eq:lyapunov} with
$\Pi_{ij}^{(r)}$ evaluated from a prescribed spectral-tensor closure. Then
the continuous evolution~\eqref{eq:cont-evolution} reproduces the
homogeneous rapid-distortion evolution of the Reynolds
stress~\eqref{eq:dRdt-rdt} exactly for that closure. With the isotropic
closure~\eqref{eq:iso-closure}, the reproduction is exact at the onset of
distortion of initially isotropic turbulence and contains no adjustable
constant, recovering the classical result~\eqref{eq:crow}.
\end{proposition}

\begin{proof}
By \eqref{eq:production-rate} the production term of
\eqref{eq:cont-evolution} contributes $\mathcal{P}_{ij}$ to
$\mathrm{d}R_{ij}/\mathrm{d}t$, and by construction~\eqref{eq:lyapunov} the
operator $B_{ij}$ contributes $\Pi_{ij}^{(r)}$. Their sum is the
right-hand side of \eqref{eq:dRdt-rdt}. For initially isotropic turbulence
$\Pi_{ij}^{(r)}$ equals the exact value~\eqref{eq:crow} by
\eqref{eq:iso-closure}--\eqref{eq:Pir-iso}.
\end{proof}

The scope of Proposition~\ref{prop:rdt} must be stated as carefully as the
result. The consistency is established at the level of the
\emph{component energies} $R_{ij}$ and is \emph{conditional on the spectral
closure}: the isotropic closure is exact only at the onset of distortion,
and its accuracy for strongly anisotropic states is that of the
isotropisation-of-production approximation it reproduces. The
\emph{spectral} distribution of the distorted energy is not enforced by
\eqref{eq:cont-evolution}; it emerges from the interaction of the strain
operator with diffusion, the cascade and the domain straining, and is
therefore a genuine prediction of the model, to be tested rather than
assumed (\S\,\ref{sec:validation}). We make no claim that the formulation
reproduces the full spectral evolution of rapid-distortion theory, which no
single-point model can; the claim is the weaker, provable one of
Proposition~\ref{prop:rdt}.

\subsection{Wavevector kinematics by domain straining}\label{sec:domain-strain}

The second ingredient of \eqref{eq:rdt}, the wavevector evolution
$\mathrm{d}k_i/\mathrm{d}t=-A_{ji}k_j$, carries the kinematic compression
and stretching of scales by the mean strain and is independent of the
amplitude evolution; in the method-of-characteristics reading of
rapid-distortion theory, the wavevectors are the characteristics along
which the amplitudes evolve. On the line, aligned with the coordinate
$x_2$, the relevant wavenumber is $k_2$. We restrict attention to the case
in which the line is aligned with a principal axis of the mean strain, so
that $A_{12}=A_{32}=0$ and a material element of the line remains a
coordinate line; this holds, in particular, for the irrotational mean flow
of a stagnation region, for which $A_{ij}$ is symmetric and its
eigenvectors define such axes. Then
\begin{equation}
  \frac{\mathrm{d}k_2}{\mathrm{d}t}=-A_{22}\,k_2 .
  \label{eq:wavenumber-line}
\end{equation}
A material element of the line of length $\ell$ stretches with the mean
flow as $\dot{\ell}/\ell=A_{22}$; since wavenumbers scale as $k_2\propto
1/\ell$, a uniform dilatation of the line reproduces
\eqref{eq:wavenumber-line} exactly,
\begin{equation}
  \frac{1}{L}\frac{\mathrm{d}L}{\mathrm{d}t}=A_{22},
  \qquad\Longrightarrow\qquad
  k_2(t)=k_2(t_0)\,\exp\!\Big(\!-\!\int_{t_0}^{t}A_{22}\,\mathrm{d}t'\Big),
  \label{eq:dilatation}
\end{equation}
where $L$ is the domain length. The dilatation~\eqref{eq:dilatation} is
implemented directly within the Lagrangian adaptive-mesh advancement of
ODT \citep{Lignell2013}, in which cell sizes already evolve, by imposing
the additional mean-flow stretching on the mesh.

This step is not redundant with the amplitude operator of
\S\S\,\ref{sec:strain-forcing}--\ref{sec:rapid-kernel}. Because the
Reynolds stress $R_{ij}$ is an intensive (per-unit-length) second moment,
it is invariant under the uniform dilatation~\eqref{eq:dilatation}: the
dilatation relabels the scales that carry the energy without changing the
component energies, while the operator $\mathcal{A}_{ij}$ changes the
component energies without relabelling scales. The two operations are the
physical-space counterparts of the characteristic motion $\boldsymbol{k}(t)$
and the amplitude evolution $\hat{u}(\boldsymbol{k}(t),t)$ in
\eqref{eq:rdt}, and together they account for the distinct contributions
of mean straining without double counting.

Two limitations are inherent to representing a three-dimensional kinematics
on one axis. The transverse wavenumbers $k_1$ and $k_3$ also strain under
$\mathrm{d}k_i/\mathrm{d}t=-A_{ji}k_j$, but are not represented on the line;
their effect on the redistribution is subsumed in the spectral closure of
\S\,\ref{sec:rapid-kernel}, which already supplies the orientation
information the line lacks. A mean flow with a rotational part, or a line
not aligned with a principal axis, tilts the material line out of the
coordinate direction and couples it to the unresolved directions; this lies
outside the present formulation and is admissible only to the extent that
the stagnation-region mean flow is irrotational.

\subsection{Summary of the formulation}\label{sec:summary}

The strain-coupled model advances the line field $u_i(y,t)$ over a
sub-step by composing three operations:
\begin{enumerate}
\item[(a)] \emph{continuous strain--diffusion phase} on a dilating domain:
  integrate $\partial_t u_i=\nu\partial_y^2 u_i+\mathcal{A}_{ij}u_j$ with
  $\mathcal{A}_{ij}=-A_{ij}+B_{ij}$, the production term exact and the
  rapid term $B_{ij}$ given by \eqref{eq:lyapunov} from the current
  Reynolds stress and the spectral closure (the baseline being
  $B_{ij}=\tfrac35 S_{ij}$, \eqref{eq:iso-B}); simultaneously dilate the
  domain by \eqref{eq:dilatation};
\item[(b)] \emph{discrete eddy events} sampled from the rate
  distribution~\eqref{eq:rate}: the triplet map~\eqref{eq:triplet} and the
  slow kernel~\eqref{eq:redistribution}--\eqref{eq:ci}, unchanged from
  standard ODT.
\end{enumerate}
Step (a) carries the mean-strain physics---production, rapid
pressure--strain and the kinematic straining of scales---and step (b)
carries the turbulence-driven physics---the inertial cascade and the slow
return to isotropy. Standard ODT is recovered exactly when $A_{ij}=0$.

The relative importance of the two steps is governed by the ratio of the
mean-strain rate $S\equiv(2S_{ij}S_{ij})^{1/2}$ to the turbulence
relaxation rate $\varepsilon/k_t$,
\begin{equation}
  \chi \equiv \frac{S\,k_t}{\varepsilon}.
  \label{eq:chi}
\end{equation}
For $\chi\gg1$ the distortion accumulates faster than the turbulence can
respond, the eddy events are negligible over the distortion time, and the
formulation reduces to the rapid-distortion limit of step (a), exact under
the closure by Proposition~\ref{prop:rdt}. For $\chi\ll1$ the turbulence
relaxes between increments of strain, the slow kernel dominates the
redistribution, and the model returns to a near-equilibrium anisotropy. The
stagnation region of a leading edge is the intermediate regime, $\chi=
O(1)$, in which neither limit is uniformly valid and both mechanisms
contribute; it is precisely there that a formulation carrying the two
together, rather than either alone, is required. The quantitative behaviour
across this range, and the fidelity of the emergent distorted spectrum, are
assessed against rapid-distortion solutions and scale-resolving data in
\S\,\ref{sec:validation}.

\bibliographystyle{jfm}
\bibliography{jfm}


\end{document}